\section{Introduction}
\label{sec:intro}
\IEEEPARstart{T}{he} Internet of Underwater Things (IoUT) links sensors, vehicles, and gateways into acoustic multi-hop meshes for environmental observation, resource exploration, and maritime awareness \cite{gupta2024uwsn,gola2023taxonomy,saleem2024ai,jiang2023orsurvey}. Acoustic links offer little bandwidth, long propagation delay, and a time-varying multipath channel \cite{gupta2024uwsn,saleem2024ai,zhu2024trustcomst}, and nodes drift with currents or are carried by vehicles, so a node's neighbourhood keeps changing \cite{gola2023taxonomy,wang2024lgdqn,liu2025mdqn,lu2025wper}. Reinforcement-learning routers have become a common response: each node learns a forwarding policy from the state of its neighbours, such as value estimates, residual energy, depth, and link quality \cite{hu2010qelar,zhang2025saqr,wang2024marv,liu2024clorP,liu2026lqarp}.

Such a router acts on \emph{held} neighbour state, not on the current state of the network, and the way that state reaches the receiver differs between routers. In some, it travels in data headers, so a node that forwards traffic advertises itself at no extra cost, while a node off the forwarding paths advertises nothing. In others, it travels only in HELLO or beacon frames, and some use both (Section~II). The two cases pose different maintenance problems. If state is carried in headers, an additional HELLO may add little. If it is carried only in HELLOs, a missing or old record can make a candidate invisible or wrongly ranked, and the packet is delayed, detoured, or lost. In both cases, every HELLO costs scarce acoustic airtime and energy, including idle listening and the preamble of a short frame.

The literature handles when to re-advertise in three ways: a fixed timer, a timer adapted to mobility \cite{li2025qlar,chen2013apu,zhu2021mobnd}, or a trigger on a void or topology change \cite{ghoreyshi2017void,li2023lrnmp}. Receive-driven replies to a newly overheard neighbour are known from adaptive beaconing in ad hoc networks \cite{chen2013apu}, and void-status beacons from void-handling protocols \cite{ghoreyshi2017void}. What remains open for underwater Q-learning routers is an evaluation that holds the routing logic fixed and varies only the maintenance policy across routers whose state travels in different ways. Most underwater studies also assert that control overhead is negligible without measuring it \cite{zhu2022roeva,zhao2023mlrs,guan2019dvor}, and most energy results leave idle listening out \cite{zhu2022roeva,zhao2023mlrs,li2023lrnmp}.

We address this with DA-HELLO, a configurable HELLO scheduler that sits beside a host router and needs only the host's neighbour query and broadcast primitive. A node sends a HELLO on a first decode of a never-recorded sender, on a change of candidate membership or void status, and, optionally, on a slow refresh; a local score-change test is also included. We evaluate it on two Q-learning routers, MARV and QELAR, and treat LQARP and SAQR as pilot hosts. Our contributions are:
\begin{enumerate}
\item \textbf{Problem.} We formulate neighbour-state maintenance for Q-learning routers as a choice that depends on how state reaches the receiver, in data headers or only in HELLOs, and pose it as minimising control overhead under a bound on the loss of delivery ratio. Without HELLOs, a HELLO-only host collapses to a delivery ratio of 0.06 to 0.18, whereas on header hosts periodic and Trickle HELLOs reduce delivery.
\item \textbf{Method.} DA-HELLO is a configurable scheduler that needs only the host's neighbour query and broadcast primitive. It sends on the first decode of a never-recorded sender, on a change of candidate membership or void status, and optionally on a slow refresh; a local score-change test is also included. We state the evaluated semantics exactly and do not claim that each HELLO is necessary for a receiver's decision.
\item \textbf{Evaluation.} Statements, margins, and seeds were fixed before the runs; comparisons are paired over 40 seeds on MARV and QELAR in header and HELLO-only form, and energy includes idle listening. On HELLO-only MARV, DA-HELLO delivers 2.1 to 3.0 percentage points more than a jittered 120\,s periodic HELLO with 43 to 53\% fewer control bytes and 8 to 10\% less energy per bit. On HELLO-only QELAR it uses 45 to 52\% fewer bytes and gains 5.0 points in the static scenario, but is 2.2 and 1.7 points lower in the drifting ones, so non-inferiority at 3 points is not shown. On header hosts it matches sending no HELLOs. The slow refresh carries delivery on HELLO-only hosts, first-contact replies harm header QELAR, and the score-change test has no demonstrated benefit. Limits of the simulation record are disclosed.
\end{enumerate}
The rest of the article is organised as follows. Section~II reviews related work and the gap. Section~III gives the preliminaries. Section~IV defines the problem and the design. Section~V presents the simulation settings and results, and Section~VI concludes.
